\documentclass[aps,prd,reprint,amsmath,amssymb,superscriptaddress,nofootinbib]{revtex4-2}
\usepackage{graphicx}
\usepackage{booktabs}
\usepackage{amsmath}
\usepackage[colorlinks=true, linkcolor=blue, citecolor=blue, urlcolor=blue]{hyperref}
\usepackage{xcolor}
\usepackage{orcidlink}
% No `array` (nor siunitx, which loads it): under revtex4-2 in this toolchain it breaks
% every p{} column. See paper/README.md.

% ---------------------------------------------------------------------------
% Every number in this paper is \input from the repository's build, never typed by
% hand; see paper/README.md. \claim{<key>} marks where each result appears, mapping it to
% its row in the repository's CLAIMS.md and giving the Lean table its section column.
\input{../../build/latex/numbers.tex}
\graphicspath{{../../build/figures/}}

\newcommand{\Nnet}{N_{\mathrm{net}}}
\newcommand{\Nkin}{N_{\mathrm{kin}}}
\newcommand{\Nins}{N_{\mathrm{ins}}}
\makeatletter
\newcommand{\claim}[1]{\expandafter\xdef\csname igc@claim@#1\endcsname{\thesection}}
\newcommand{\claimsection}[1]{\@ifundefined{igc@claim@#1}{--}{\csname igc@claim@#1\endcsname}}
\makeatother

\begin{document}

\title{When induced gravity attracts: an exact count for the Standard Model}

\author{David T. Swanson~\orcidlink{0009-0003-0580-6476}}
\noaffiliation

\date{\today}

\begin{abstract}
In induced gravity, Newton's constant arises from the one-loop vacuum energy of
matter, and its sign is the sign of a sum over species, $\Nnet=\sum\sigma(N-6c)$,
with $\sigma$ the statistics sign, $N$ the number of field components and $cR$ the
trace of the curvature term in the species' own fluctuation operator. Under a
spectral cutoff taken in each species' own operator, $\Nnet$ is \smNet{} for the
Standard Model with the Higgs minimally coupled, in units of one minimally coupled
scalar, the value an earlier species count also reaches. Its kinetic parts $\sigma N$
sum to \smKinetic{} and its curvature insertions $-6\sigma c$ to \smInsertion{}. The
cutoff weights each insertion at exactly six times its kinetic part, for every cutoff
function of bounded variation, so the count is exact and does not depend on the cutoff
function; the sign would reverse only at a weighting of \smThreshold{}, \smMarginPct{}
below six, which only a cutoff departing from this condition could reach. The chain
from the regulated action to the sign is proved for proper-time cutoffs, and for other
cutoff functions rests on the spectral-action expansion as the literature states it. A
horizon-entropy count, entanglement \smEntanglement{} against the gauge fields' contact
term \smContact{}, agrees with it species by species, a consistency check on the
species table. The count is an integer and \smNet{} is its least positive value; with
the Standard Model's gauge group, three generations and one Higgs doublet are the only
content of whole generations and doublets that reaches it. Attraction begins at
\floorGenerationsSpelled{} generations, a floor rather than a selection. The count
requires the Higgs doublet's common curvature coupling to stay below
\smCouplingBound{}, and keeps its sign for every number of right-handed neutrinos. If
Newton's constant is entirely induced, as Sakharov proposed, the count becomes a
condition on content, much as anomaly cancellation is: each gauge boson below the cutoff
must be paid for by \matterPerBosonSpelled{} Weyl fermions or minimally coupled real
scalars, and a conformally coupled Higgs doublet is excluded, while Higgs inflation,
whose coupling has the other sign, is not. The counts and the identities they rest on
are machine-checked in Lean.
\end{abstract}

\maketitle

% =============================================================================
\section{Introduction}
\label{sec:intro}

Sakharov proposed that the Einstein action is not fundamental: it is the change in
the one-loop action of quantum fields when spacetime is curved~\cite{Sakharov1967}.
Amati and Veneziano extended the idea to the gauge couplings, inducing them from the
same matter content~\cite{AmatiVeneziano1981}, and Visser reviews the mechanism in
modern form~\cite{Visser2002}. In each version the coefficient of $\int\!\sqrt{g}\,R$
in the induced action is a sum over species, each weighted by the first
heat-kernel coefficient of its own operator~\cite{BirrellDavies1982,Vassilevich2003}.

That makes the \emph{sign} of the induced Newton constant the sign of a sum over
species, the net count $\Nnet$ of Eq.~\eqref{eq:net}. Each species contributes a
kinetic part, from the Laplacian acting on its field components, and a curvature
insertion, from the curvature term in its own operator. For the fermions and the
gauge fields the two have opposite signs (for minimally coupled scalars and for
ghosts the insertion vanishes), so whether the matter we have induces attractive
gravity depends on what that matter is. In units of one minimally coupled scalar, a
scalar counts \weightScalar{}, a Weyl fermion \weightWeyl{}, and a gauge field with
its two ghosts \weightGaugeField{}. The Standard Model's \smScalars{} scalars and
\smWeyl{} Weyl fermions give \smNetWithoutGauge{}, its \gaugeVectors{} gauge fields
\smGaugeSector{}, and the net is \smNet{}. Figure~\ref{fig:generations} shows the
count for the Standard Model's gauge group, one Higgs doublet, and any number of
generations.

\begin{figure}[t]
\includegraphics[width=\columnwidth]{generations.pdf}
\caption{The net count $\Nnet$ of Eq.~\eqref{eq:net}, in units of one minimally
coupled scalar, against the number of generations, with the Standard Model's gauge
group and one Higgs doublet. Gravity attracts exactly when $\Nnet>0$. The count is
negative up to two generations and positive from three on, where it is
\genCountThree{} (marked, since the bar is too short to see).}
\label{fig:generations}
\end{figure}

The count has been made before. Tanaka computed the induced Einstein action in
supersymmetric models, found that vector multiplets induce a negative Newton constant
and scalar multiplets a positive one, and concluded that the minimal supersymmetric
Standard Model and minimal supersymmetric SU(5) need three generations or more for
attractive gravity~\cite{Tanaka1996}. Broda and Szanecki weighted the species exactly
as here, under a Planckian proper-time cutoff~\cite{BrodaSzanecki2009}. For the
Standard Model with its Higgs doublet and twelve gauge fields, their combination of
species equals one, and they note that the black hole's geometric entropy carries the
same combination. They call the gauge fields' contribution disputable, since the
gauge fields are themselves induced. The net count of this paper is therefore not
new.

What this paper adds is the count's structure, and its conditions. Resolved into its
parts, the kinetic parts sum to \smKinetic{} and the curvature insertions to
\smInsertion{}. The weighting that the cutoff gives the insertions is exactly six for
every profile of bounded variation, by an identity between the profile's moments, so
the count is exact and the profile cannot change it; for proper-time profiles, the
heat kernel's among them, the whole chain from the regulated action to the sign is
proved, the one-to-six ratio derived rather than quoted (Sec.~\ref{sec:weights},
Appendices~\ref{app:setup}--\ref{app:onetosix}). An exact criterion locates the
boundary: the sign reverses at a weighting of \smThreshold, \smMarginPct{} below six,
which only a cutoff departing from that condition could reach (Sec.~\ref{sec:count}).
A horizon-entropy count agrees with the count species by species, a consistency check
on the species table and on condition K2 (Sec.~\ref{sec:horizon}). The count is an
integer, so \smNet{} is the least positive value it can take, and with the gauge group
given, the Standard Model's generations and doublet are the only content of whole
generations and Higgs doublets that reaches it. The same count gives a floor on the
number of generations, a bound on the Higgs doublet's curvature coupling, and the
effect of right-handed neutrinos (Sec.~\ref{sec:corollaries}). Under Sakharov's
premise that Newton's constant is entirely induced, it becomes a condition on content:
gauge bosons must be paid for in matter, and a conformally coupled Higgs doublet is
excluded (Sec.~\ref{sec:forbids}).
Every count, and the identities behind them, is machine-checked in Lean and
cross-checked against a second implementation, both reading one table of species
with a citation for each entry (Sec.~\ref{sec:lean}).

The paper does not claim the magnitude of Newton's constant, which moves with the
cutoff's scale and the profile's shape and is reported only as a range
(Sec.~\ref{sec:magnitude}). Nor does it claim anything about an induced graviton's
dynamics, or about the vacuum energy (Sec.~\ref{sec:unsettled}).

% =============================================================================
\section{The conditions}
\label{sec:conditions}

The count is made under four conditions, stated in Table~\ref{tab:conditions}. Of
the four, K1 is the one the sign depends on; K2 decides whether the count is
\smNetWithoutGauge{} or \smNet; K3 and K4 limit what is claimed.

\begin{table*}
\caption{The conditions the count is made under.}
\label{tab:conditions}
\centering
\input{../../build/latex/tables/conditions.tex}
\end{table*}

Under K1, with the cutoff taken in each species' own fluctuation operator, the
curvature term of that operator enters the heat kernel's first coefficient on the
same footing as the Laplacian, and the insertion carries the same moment of the
cutoff function as the kinetic part (Sec.~\ref{sec:weights}). A cutoff placed in
another operator can reweight the insertion, and the criterion of
Sec.~\ref{sec:count} then decides whether the sign survives.

Condition K2 takes the gauge fields as fundamental fields up to the cutoff, each with
its two ghosts and its contact term with the horizon. It is the condition Broda and
Szanecki question, since the gauge fields may themselves be
induced~\cite{BrodaSzanecki2009}. Without the gauge fields the Standard Model's count
would be \smNetWithoutGauge{}; counting them with their ghosts adds \smGaugeSector{}
and brings it to \smNet{}. K2 therefore leaves the sign positive and brings the count
to its least positive value.

The regulator here is covariant. Condition K3 records that what is computed is the
coefficient of $\int\!\sqrt g\,R$ in the Euclidean action on a fixed background, not
an induced graviton's propagation, whose time and space derivative terms a
non-covariant regulator could treat differently~\cite{Visser2002}. Condition K4 sets
the cosmological term aside.

\paragraph{Further assumptions.}
Besides K1 to K4, the count assumes: free massless fields on a purely gravitational
background; a closed Euclidean background with $\int\!\sqrt g\,R\neq0$, which makes the
induced coefficient unique; a cutoff function whose moment $M_f$ (Eq.~\eqref{eq:moment})
is positive, which holds for every proper-time cutoff and every non-increasing one;
each gauge vector in Feynman gauge with its two Faddeev--Popov ghosts, counted as an
abelian field, which is exact at one loop on a purely gravitational background, with
the dependence on the gauge parameter not examined; and, for the main result, a
minimally coupled Higgs. Appendix~\ref{app:setup} states each in the form the proofs
use.

% =============================================================================
\section{Each species' weight}
\label{sec:weights}

A \emph{species} is a field with its own fluctuation operator; the Faddeev--Popov
ghosts count as species of their own, and a \emph{content} is a list of species with
their multiplicities. Each species' fluctuation operator is of Laplace type,
$D=D_0+X$ with $D_0=-\nabla^2$ acting on $N$ field components, and a curvature term
$X$ whose trace is $\operatorname{tr}X=cR$. Its heat trace has the
expansion~\cite{Vassilevich2003}
\begin{equation}
\operatorname{Tr}e^{-\tau D}=\tau^{-d/2}\bigl[a_0+a_1\tau+O(\tau^2)\bigr],
\label{eq:heat}
\end{equation}
whose curvature coefficient is
\begin{equation}
a_1=(4\pi)^{-d/2}\Bigl(\frac{N}{6}-c\Bigr)\!\int\!\sqrt{g}\,R .
\label{eq:a1}
\end{equation}
The Euclidean one-loop action of a content is
$W=\tfrac12\sum_i\sigma_i\log\det D_i$, with $\sigma_i$ the statistics sign. The
cutoff of K1 regulates the logarithm species by species: $\log\det D$ is replaced by
$-\operatorname{Tr}f(D/\Lambda^2)$, for a cutoff function $f$ (the \emph{profile}),
such as the sharp step $\theta(1-u)$, the heat kernel's own $e^{-u}$, or a Gaussian
$e^{-u^2}$. So $W_f(\Lambda)=-\tfrac12\sum_i\sigma_i\operatorname{Tr}f(D_i/\Lambda^2)$.
At order $\Lambda^{d-2}$ each trace carries the heat kernel's $a_1$ times one moment
of the profile~\cite{ChamseddineConnes1997,vanSuijlekom2024,EcksteinIochum2018,
EstradaGraciaBondiaVarilly1998},
\begin{equation}
M_f=\frac{1}{\Gamma(s)}\int_0^\infty f(u)\,u^{s-1}\,du ,\qquad s=\frac d2-1 ,
\label{eq:moment}
\end{equation}
normalized to the heat kernel's own profile, for which $M_{e^{-u}}=1$. In the
Euclidean convention $W\supset-(16\pi G)^{-1}\!\int\!\sqrt{g}\,R$~\cite{FrolovFursaevZelnikov1997},
the induced coupling is then
\begin{equation}
\frac{1}{16\pi G}=\frac{M_f}{12\,(4\pi)^{d/2}}\,\Nnet\,\Lambda^{d-2},
\label{eq:newton}
\end{equation}
with the content's net count
\begin{equation}
\Nnet=\sum_i\sigma_i\,(N_i-6c_i),
\label{eq:net}
\end{equation}
so the sign of $G$ is the sign of $\Nnet$ for every profile with a positive moment,
and gravity attracts exactly when $\Nnet>0$.

\claim{species-weight}%
Each species' weight $\sigma(N-6c)$ is the sum of a \emph{kinetic part} $\sigma N$,
from $D_0$, and a \emph{curvature insertion} $-6\sigma c$, from $X$.
Table~\ref{tab:species} gives the four species the count needs. The weights reproduce
the cited values: \weightScalar{} for a minimal scalar, \weightWeyl{} for a Weyl
fermion, and \weightGaugeField{} for a gauge field with its two
ghosts~\cite{BrodaSzanecki2009,Kabat1995}. The conventions are in
Appendix~\ref{app:counts}.

\begin{table*}
\caption{Each species' statistics sign $\sigma$, field components $N$, curvature trace
$c$ and weight, with its coefficient in a horizon's entanglement entropy and its
contact term with the horizon (Sec.~\ref{sec:horizon}), both in units of one minimally
coupled real scalar. The entanglement coefficient is not a count of physical degrees
of freedom: a Weyl fermion has two, which count half each, as fermions do. The
contact term is Kabat's~\cite{Kabat1995}.}
\label{tab:species}
\centering
\input{../../build/latex/tables/species.tex}
\end{table*}

The claim that the profile cannot change the sign rests on two steps, with different
scopes.

\paragraph{Step 1: the moments.}
\claim{one-to-six}\claim{profile-family-weighting}%
The insertion can be read in two ways. In Eq.~\eqref{eq:a1} it is part of the full
operator's coefficient, and carries the profile's moment $M_f$. Read instead as the
first-order response of the cutoff trace to the curvature term,
$\frac{d}{d\varepsilon}\operatorname{Tr}f\bigl((D_0+\varepsilon X)/\Lambda^2\bigr)$ at
$\varepsilon=0$, as a perturbative treatment of $X$ would compute it, it carries the
profile's \emph{derivative} moment, $-\int_0^\infty f'(u)\,u^{s}du$, divided by
$\Gamma(s+1)$. The two agree because of an integration by parts,
\begin{equation}
-\int_0^\infty u^{s}\,df(u)=s\int_0^\infty f(u)\,u^{s-1}du ,
\label{eq:parts}
\end{equation}
which holds for every profile whose derivative is a measure: every profile of
bounded variation with finite moments, the sharp step included. For the sharp step,
$f'=-\delta(u-1)$, the left side is one and the right side is $s\cdot(1/s)$. So the
insertion's weight relative to the kinetic part is six for every profile of bounded
variation (Appendix~\ref{app:onetosix}). This step is an identity between moments; it
does not yet say that the moments are what the cutoff trace carries.

\paragraph{Step 2: the expansion.}
\claim{sign-profile-independent}%
That the cutoff trace carries $M_f\,a_1$ at order $\Lambda^{d-2}$, so that the profile
enters Eq.~\eqref{eq:newton} as one positive factor, is the large-$\Lambda$ expansion of
$\operatorname{Tr}f(D/\Lambda^2)$ in heat-kernel coefficients. Its standing depends on
the profile:
\begin{itemize}
\item For \emph{proper-time} profiles, $f(u)=\int e^{-tu}\,d\nu(t)$ with $\nu$ a finite
measure on a bounded proper-time interval, the heat kernel's own profile among them,
the expansion is a theorem: it follows from the heat trace alone
(Appendix~\ref{app:abelian}). For them the two readings of the insertion are both
computed and shown to agree, so the one-to-six ratio is derived, not quoted
(Appendix~\ref{app:onetosix}), for an insertion that is constant over the background.
\item For other \emph{Laplace-transform} profiles and for \emph{Schwartz} profiles the
expansion is stated in the spectral-action
literature~\cite{vanSuijlekom2024,EcksteinIochum2018,EstradaGraciaBondiaVarilly1998}.
It is taken as a hypothesis for them, not proved. The spectral action of
Ref.~\cite{ChamseddineConnes1997} likewise puts its cutoff in an operator, the Dirac
operator, though there the cutoff trace defines the bosonic action rather than
regulating each species' one-loop trace.
\item For the \emph{sharp step} the identity~\eqref{eq:parts} holds exactly, but its
cutoff trace has the expansion beyond the leading term only as a Ces\`aro
mean~\cite{EstradaGraciaBondiaVarilly1998}: the subleading $\Lambda^{d-2}$ term exists
only after averaging over the cutoff scale. The sharp step is therefore covered by
Step~1, and by Step~2 only in that averaged sense.
\end{itemize}

% =============================================================================
\section{The Standard Model's count}
\label{sec:count}

\paragraph{The count.}
\claim{headline-count}\claim{standard-model-count}\claim{induced-newton-constant-positive}%
The Standard Model with one Higgs doublet carries \smScalars{} real scalars,
\smWeyl{} Weyl fermions in \smGenerations{} generations, \gaugeVectors{} gauge
vectors and their \smGhosts{} ghosts~\cite{PDG2024}. Table~\ref{tab:count} gives the
count species by species. The kinetic parts sum to $\Nkin=\smKinetic$ and the
insertions to $\Nins=\smInsertion$, a net of \smNet{}: with the Higgs minimally
coupled, the Standard Model induces a positive Newton constant under K1.

\begin{table*}
\caption{The Standard Model's count, species by species, in units of one minimally
coupled scalar: each species' kinetic part is $m\sigma N$ and its insertion
$-6m\sigma c$, for multiplicity $m$. The last two columns are the horizon count's
(Sec.~\ref{sec:horizon}).}
\label{tab:count}
\centering
\input{../../build/latex/tables/count.tex}
\end{table*}

\paragraph{The criterion.}
\claim{criterion}\claim{criterion-margin}%
Weight every species' insertion by $\rho$ in place of six. The count
$\Nnet(\rho)=\Nkin+(\rho/6)\Nins$ is positive exactly when
\begin{equation}
\rho>\frac{6|\Nkin|}{\Nins}=\smThreshold\approx\smThresholdDecimal .
\label{eq:criterion}
\end{equation}
Condition K1 gives $\rho=6$ for every profile of bounded variation, so under K1 $\rho$
is not a free parameter and the count is exactly \smNet. The criterion locates the
boundary instead. The \emph{margin}, the threshold's distance below six as a fraction
of six, is $(6-\smThreshold)/6=\smMargin$, or \smMarginPct{}; it equals
$\Nnet/\Nins$, the net count over the insertions' total, and is \nuTwoMargin{} and
\nuThreeMargin{} with two and three right-handed neutrinos (Table~\ref{tab:neutrinos}).
It is how far a scheme departing from K1 would have to lower the insertions' common
weight to reverse the sign; no such scheme is analysed here. The criterion takes one
$\rho$ for every species: a scheme that
reweighted fermion and vector insertions differently, whose contributions have
opposite signs, would need the criterion species by species.
Figure~\ref{fig:criterion} shows the count against $\rho$ for the Standard Model with
no, two and three right-handed neutrinos (Sec.~\ref{sec:corollaries}).

\begin{figure}[t]
\includegraphics[width=\columnwidth]{criterion.pdf}
\caption{The net count against $\rho$, the weight of every species' curvature
insertion relative to its kinetic part (Eq.~\eqref{eq:criterion}), for the Standard
Model with no, two and three right-handed neutrinos. Each line crosses zero at its
content's threshold (dots): \smThreshold, \nuTwoThreshold{} and \nuThreeThreshold{}
respectively. The cutoff of K1 gives $\rho=6$ (dashed).}
\label{fig:criterion}
\end{figure}

% =============================================================================
\section{The horizon route}
\label{sec:horizon}

\claim{horizon-independent-route}\claim{horizon-count}%
Larsen and Wilczek showed that the one-loop correction to a black hole's entropy and
the renormalization of Newton's constant carry the same
divergence~\cite{LarsenWilczek1996}. In induced gravity, the horizon entropy is the
fields' entanglement entropy less contact terms with the
horizon~\cite{FrolovFursaevZelnikov1997}, and once $G$ is the induced constant, the
Bekenstein--Hawking entropy $A/4G$ requires the two counts to agree. Each species
enters the entanglement entropy with a coefficient that, in units of one real scalar,
is $1$ for a scalar, $1$ for a Weyl fermion (whose two physical degrees of freedom
count half each, as fermions do), $2$ for a gauge field's two polarizations, and $0$
for a ghost (Table~\ref{tab:species}). A gauge field adds a contact term with the
horizon, \contactPerGaugeField{} in units of one minimal scalar~\cite{Kabat1995},
which Donnelly and Wall read as the entanglement of edge
modes~\cite{DonnellyWall2015}; the count uses only its coefficient, which the two
readings share. Solodukhin reviews both relations~\cite{Solodukhin2011}.

The last two columns of Table~\ref{tab:count} give the Standard Model's horizon count:
entanglement \smEntanglement{} and contact term \smContact{}, a net of \smNet{}, as
the proportionality requires. The agreement holds species by species. A gauge field
costs four by either route: in the induced count its weight \weightVector{} and its two
ghosts' \weightGhost{} each, and in the horizon count its two polarizations and its
contact term \contactPerGaugeField. The agreement therefore holds for every content
whose vectors each carry their two ghosts. It is a
consistency check on the species table and on condition K2, not independent evidence
for the sign of the Standard Model's count. Broda and Szanecki note the same
agreement for the geometric entropy, their name for the entanglement
entropy~\cite{BrodaSzanecki2009}.

% =============================================================================
\section{Corollaries}
\label{sec:corollaries}

A content with the Standard Model's gauge group, $g$ generations, $h$ Higgs doublets
and $n$ right-handed neutrinos has the count
\begin{equation}
\Nnet=\weylPerGeneration\,g+\scalarsPerDoublet\,h+n\smGaugeSector ,
\label{eq:content}
\end{equation}
the gauge sector contributing \smGaugeSector{}. The floor, the least content and the
right-handed neutrinos' effect hold for every $g$, $h$ and $n$; the coupling bound is
stated for the Standard Model, and in general form in Appendix~\ref{app:counts}.

\paragraph{A floor on generations.}
\claim{generation-floor}%
With one doublet, the count runs \genCountZero, \genCountOne, \genCountTwo{} and
\genCountThree{} for zero to three generations (Fig.~\ref{fig:generations}), and it is
positive for every number of generations from
\floorGenerationsSpelled{} on. Attraction begins at \floorGenerationsSpelled{}
generations. That is a floor, not a selection: every larger number also attracts.
Tanaka found the same floor for supersymmetric contents~\cite{Tanaka1996}; here it
holds for the Standard Model's own content, with the gauge fields counted, and three
generations clear it by one.

\paragraph{The least content.}
\claim{unique-least-content}%
The count is an integer, so \smNet{} is the least positive value it can take. Among
contents of whole generations and Higgs doublets with this gauge group, the only one
that reaches it is \countOneGenerationsSpelled{} generations and
\countOneDoubletsSpelled{} doublet. By Eq.~\eqref{eq:content} a count of one
needs $\weylPerGeneration\,g+\scalarsPerDoublet\,h=\countOneNeeds$, which allows at
most three generations, and of those only three with one doublet solves it. With a
right-handed neutrino in every generation, every term of the count is a multiple of
four, so no such content has a count of one. Both statements hold for every number of
generations and doublets; Appendix~\ref{app:counts} gives the finite search that
cross-checks them. Other representations are outside the statement: adding singlet
scalars, for instance, would give other contents of count one.

\paragraph{The Higgs coupling.}
\claim{coupling-bound}%
Give the doublet's \smScalars{} real scalars a common coupling $\xi$ to curvature,
with $\xi=0$ minimal and $\xi=1/6$ conformal~\cite{BirrellDavies1982}. Each scalar's
curvature trace becomes $\xi$, and the count falls by $6\xi$ per scalar. Attraction
then requires
\begin{equation}
\xi<\smCouplingBound ,
\end{equation}
strictly, since the count vanishes at the bound. With two and three right-handed
neutrinos the bound rises to \nuTwoCouplingBound{} and \nuThreeCouplingBound. A
coupling below the bound, of either sign, keeps the count positive; what the bound
excludes, if gravity is induced, is taken up in Sec.~\ref{sec:forbids}.

\paragraph{Right-handed neutrinos.}
\claim{neutrino-margins}%
\claim{neutrino-sign-robust}%
Each right-handed neutrino adds one to the count. With two the count is \nuTwoNet,
the threshold \nuTwoThreshold{} and the margin \nuTwoMarginPct; with three,
\nuThreeNet, \nuThreeThreshold{} and \nuThreeMarginPct{} (Table~\ref{tab:neutrinos}).
The rows are none, the Standard Model's own content; two, the least number that gives
the two non-zero neutrino masses the observed oscillations need~\cite{PDG2024}, as in
the minimal seesaw of Frampton, Glashow and Yanagida~\cite{FramptonGlashowYanagida2002};
and three, one per generation. Eq.~\eqref{eq:content} gives the count for any number,
and the sign holds for every one.

\begin{table}
\caption{The Standard Model with $\nu_R$ right-handed neutrinos added: the count, the
threshold weighting, the margin, and the bound on the Higgs doublet's common curvature
coupling.}
\label{tab:neutrinos}
\centering
\input{../../build/latex/tables/neutrinos.tex}
\end{table}

% =============================================================================
\section{The magnitude, and why it is not claimed}
\label{sec:magnitude}

\claim{magnitude-range}%
In four dimensions Eq.~\eqref{eq:newton} gives the cutoff at which a content induces
Newton's constant,
\begin{equation}
\Lambda=M_P\sqrt{\frac{12\pi}{M_f\,\Nnet}} ,\qquad M_P=G^{-1/2}.
\end{equation}
The profile's moment is one for the sharp step and for the heat kernel, and
\momentGaussian{} for the Gaussian $e^{-u^2}$ (Appendix~\ref{app:magnitude}). For the
Standard Model the required cutoff ranges from \smCutoffLow{} to
\smCutoffHigh{}\,$M_P$ over these profiles (Table~\ref{tab:magnitude}), and
right-handed neutrinos lower it. That a net of one needs a cutoff a few times the
Planck mass is the same statement as Broda and Szanecki's coefficient $1/12\pi$ at a
Planckian cutoff~\cite{BrodaSzanecki2009}. The paper does not predict the cutoff's
scale or its profile; given them, this equation fixes the magnitude. The magnitude
moves with the profile. The sign does not.

\begin{table}
\caption{The profile's moment $M_f$ and the cutoff in Planck units, by profile and
number of right-handed neutrinos $\nu_R$. The heat kernel's row rests on the proved
expansion, the Gaussian's on the imported one for Schwartz profiles; the sharp step's
moment is exact, but its expansion beyond the leading term holds only as a Ces\`aro
mean (Appendix~\ref{app:magnitude}).}
\label{tab:magnitude}
\centering
\input{../../build/latex/tables/magnitude.tex}
\end{table}

% =============================================================================
\section{What the count forbids, if gravity is induced}
\label{sec:forbids}

Everything so far concerns the sign of the curvature term a content induces. It becomes
a constraint on content under one further premise: that Newton's constant is
\emph{entirely} induced, with no bare term, by the species lighter than the cutoff. That
is Sakharov's premise~\cite{Sakharov1967}, and this paper does not argue for it. Without
it, a bare term can carry the sign, and the count gives only the sign of matter's
contribution to $1/G$; nothing below is then forbidden. With it, and under K1 to K4, a
content whose count is not positive induces no attractive gravity. The count then acts
as a condition on what the content can be, much as anomaly cancellation does: it
excludes contents without predicting a number. Unlike anomaly cancellation, it holds
only under the premise.

\paragraph{Gauge bosons must be paid for.}
\claim{gauge-boson-cost}\claim{extension-rule}%
A gauge boson with its two ghosts contributes \weightGaugeField{} by either route
(Sec.~\ref{sec:horizon}), and a Weyl fermion or a minimally coupled real scalar
contributes \weightScalar. A content of these species therefore has
$\Nnet=n_{\text{scalar}}+n_{\text{Weyl}}-\matterPerBoson\,n_{\text{vector}}$, the
combination of Broda and Szanecki~\cite{BrodaSzanecki2009}, and the Standard Model,
with a count of \smNet, has one unit to spare. Hence:
\begin{itemize}
\item the Standard Model with one further gauge boson below the cutoff, and nothing
else, has count \smPlusOneBoson;
\item with a complex scalar added, minimally coupled, to give that boson its mass, the
count is \smPlusBosonAndHiggs;
\item the Standard Model extended by $b$ gauge bosons and $k$ further Weyl fermions or
minimally coupled real scalars has a count $k-\matterPerBoson\,b$ above the Standard
Model's, and attracts exactly when $k\ge\matterPerBoson\,b$
(Appendix~\ref{app:counts}).
\end{itemize}
Each new gauge boson must bring \matterPerBosonSpelled{} such fields with it. Under the
premise this bears on every extension with a new gauge boson, such as a dark photon or
a $Z'$: its content below the cutoff must pay for its bosons. A boson's mass does not
change what it costs. A mass $\mu$ shifts each species' operator by the constant
$\mu^2$, which changes $a_1$ only by $-\mu^2a_0$ (Appendix~\ref{app:onetosix}), a
multiple of the volume rather than of $\int\!\sqrt g\,R$; so at order $\Lambda^{d-2}$ a
species lighter than the cutoff counts as a massless one does. The formal statements are
for massless fields (Sec.~\ref{sec:unsettled}) and do not include this step.

\paragraph{The Higgs doublet's curvature coupling.}
\claim{conformal-higgs-excluded}%
Under the premise the doublet's common coupling must satisfy $\xi<\smCouplingBound$
(Sec.~\ref{sec:corollaries}). The conformal coupling lies above the bound: with the
doublet conformally coupled, $\xi=1/6$, the Standard Model's count is
\smConformalCount, so a conformally coupled Higgs doublet is excluded. Right-handed
neutrinos leave this in place: with two the bound is \nuTwoCouplingBound, and with three
it is \nuThreeCouplingBound, the conformal value itself, where the count vanishes.
Higgs inflation is not excluded. It needs a large coupling, $\xi\approx49000\sqrt\lambda$
in the convention of Ref.~\cite{BezrukovShaposhnikov2008}, which places the conformal
coupling at $-1/6$. In this paper's convention the same coupling is large and negative,
and it raises the count, by \smCouplingSlope{} per unit of $|\xi|$, rather than lowering
it. The bound excludes large couplings of the conformal sign, not Higgs inflation's.

\paragraph{What is not forbidden.}
Further generations, right-handed neutrinos and minimally coupled Higgs doublets all
raise the count, so it allows them. The floor of \floorGenerationsSpelled{} generations
excludes only what observation already excludes.

\paragraph{How the constraints meet observation.}
Both constraints concern the complete content below the cutoff, which lies at a few
Planck masses (Sec.~\ref{sec:magnitude}). The matter that pays for a gauge boson may lie
anywhere below it, so no single discovery can violate the gauge-boson rule: a new gauge
boson found without the matter to pay for it would only place that matter higher. The
rule constrains complete models, which must list their content up to the cutoff. The
coupling bound is likewise a constraint on models rather than a prediction for a
measurement: collider data bound the Higgs boson's curvature coupling only at the order
of $10^{15}$~\cite{AtkinsCalmet2013}, far above $\smCouplingBound$.

% =============================================================================
\section{What the count does not settle}
\label{sec:unsettled}

\paragraph{Whether gravity is induced.} The constraints of Sec.~\ref{sec:forbids} rest
on Sakharov's premise that Newton's constant is entirely induced, which the paper
states and does not argue for.

\paragraph{The dynamical sector.} The count is the coefficient of
$\int\!\sqrt g\,R$ on a fixed Euclidean background, not an induced graviton's
propagation; with a regulator that is not Lorentz-invariant the two need not
agree~\cite{Visser2002} (K3).

\paragraph{The vacuum energy,} the cutoff trace's leading term, is set aside (K4).

\paragraph{Beyond one loop.} The count is the one-loop coefficient of massless fields
on a purely gravitational background. Interactions, masses and higher loops are not
included.

\paragraph{Other cutoffs.} A cutoff taken in an operator other than the species'
own fluctuation operator falls outside K1. Such a cutoff can move part of the
curvature term out of the regulated operator, so the insertion's weight need no
longer be six, and the count is not made for it here. The sharp step is covered
only in the Ces\`aro sense noted in Sec.~\ref{sec:weights}.

\paragraph{The content.} The main result takes the Higgs minimally coupled, and the
corollaries range over whole generations, Higgs doublets and right-handed neutrinos,
not over other representations.

% =============================================================================
\section{Machine-checked statements}
\label{sec:lean}

The counts and the identities above are stated and proved in Lean~4 with Mathlib.
The Lean proofs use a species table generated from a data file; a second
implementation, in Python, reads the same file, and the build fails if the two ever
disagree on a count. The proofs contain no unproved steps (Lean's \texttt{sorry}) and
no unverified evaluation (\texttt{native\_decide}), and use no axioms beyond Lean's
standard three: propositional extensionality, choice and quotient soundness.

What is proved: the weight's decomposition into kinetic part and insertion; the
integration by parts~\eqref{eq:parts} for every profile of bounded variation; the
expansion of the cutoff trace for proper-time profiles, from the heat trace alone;
the one-to-six identity, with the insertion's first-order response computed; the
induced coupling~\eqref{eq:newton} of any content; the Standard Model's and the
horizon counts; the criterion; the floor and the least content, for every number of
generations and doublets; the coupling bound, for every real $\xi$; the effect of
right-handed neutrinos, for any number of them; the magnitude formula with the three
profiles' moments; and the cost of a gauge boson by both routes, the count of every
content of scalars, Weyl fermions and gauge vectors, and the Standard Model's
extensions and its conformally coupled Higgs doublet, with the Newton constant each
induces.

What is imported, as hypotheses the theorems take rather than as axioms: the heat
trace's expansion and its first coefficients~\cite{Vassilevich2003}; the cutoff
trace's expansion for profiles outside the proper-time class~\cite{vanSuijlekom2024,
EcksteinIochum2018,EstradaGraciaBondiaVarilly1998}; and the species table's entries,
including the contact term~\cite{Kabat1995}. The sign convention for the induced
coupling is a definition~\cite{FrolovFursaevZelnikov1997}.

\paragraph{What the checks do not reach.} A proof shows that a statement holds; it
does not show that the statement says what the result citing it says, that the
imported inputs and the species table's entries are the right physics, or that the
conditions are the right ones. Table~\ref{tab:lean} in Appendix~\ref{app:formal} sets
each result beside the section it appears in and the statements it cites. The two
implementations, written by one
author from one table, guard against slips of arithmetic rather than against a shared
misreading.

% =============================================================================
\section{Conclusion}
\label{sec:conclusion}

Gravity's attractive sign can be read off the matter we have. Under a spectral cutoff
in each species' own operator, the Standard Model's count is exactly \smNet{}, the
least positive value the count can take. Its kinetic parts, \smKinetic, and its
insertions, \smInsertion, stand at one to six for every profile of bounded variation,
so the profile cannot change the count (the chain to the sign is proved for
proper-time profiles, and rests on the stated expansion for others); the sign would
reverse only under a cutoff that lowered the insertions' weight by \smMarginPct. The
count agrees species by species with the horizon count and keeps its sign for every
number of right-handed neutrinos. Attraction begins at \floorGenerationsSpelled{}
generations, a floor rather than a selection, and requires the Higgs doublet's
curvature coupling to stay below \smCouplingBound{}: any model that induces gravity
from this content under K1 to K4 must meet both. If Newton's constant is entirely
induced, the count becomes a condition on content: each gauge boson must be paid for by
\matterPerBosonSpelled{} Weyl fermions or minimally coupled real scalars, and a
conformally coupled Higgs doublet is excluded.

\paragraph{Data and code availability.}
% The archive's DOI is the SwansonIGC entry's, in data/references.yaml.
Every number, table and figure in this paper is generated by the
\texttt{induced-gravity-count} repository~\cite{SwansonIGC}, where one command reproduces them from a
clean checkout and the Lean proofs build with the paper's numbers. Its
\texttt{CLAIMS.md} maps each result to the computation and the Lean statements behind
it, and its \texttt{LIMITS.md} states what is not established.

\appendix

%======================================================================
\section{Setup: the regulated action and the Newton constant it induces}
\label{app:setup}

\paragraph{Spectral data.}
Each species' fluctuation operator $D$ enters only through its spectral measure $\mu$,
which counts eigenvalues with multiplicity. Its heat trace, and its cutoff trace under a
profile $f$ at scale $L=\Lambda^2$, are
\begin{equation}
\begin{gathered}
  \Theta(\tau)=\int e^{-\tau\lambda}\,d\mu(\lambda),\\
  \operatorname{Tr} f(D/L)=\int f(\lambda/L)\,d\mu(\lambda).
\end{gathered}
\end{equation}
For every species, $\Theta(\tau)$ is assumed to exist for all $\tau>0$ and, with
$s=d/2-1$, to satisfy
\begin{equation}\label{eq:heatbound}
  \bigl|\tau^{s+1}\Theta(\tau)-a_0-a_1\tau\bigr|\le C\tau^2\qquad\text{for } 0<\tau\le\tau_0,
\end{equation}
the small-$\tau$ expansion of Ref.~\cite{Vassilevich2003}. In four dimensions $s=1$.

\paragraph{The background.}
The background is closed, with $\int\!\sqrt g\,R\neq0$, and a species of $N$ components
whose operator carries the curvature term $cR$ has
$a_1=(4\pi)^{-(s+1)}\,(N/6-c)\int\!\sqrt g\,R$~\cite{Vassilevich2003,BirrellDavies1982}.
The non-vanishing integral makes the induced coefficient below unique.

\paragraph{The coefficient at order $L^s$ (condition K4).}
A function $T(L)$ has coefficient $b$ at order $L^s$ when, for some $A$,
\begin{equation}\label{eq:coeff}
  L^{-s}\bigl(T(L)-A\,L^{s+1}\bigr)\to b\qquad (L\to\infty).
\end{equation}
The leading term $A\,L^{s+1}$ is the cosmological one. It is removed by the definition
and never claimed. The coefficient, when it exists, is unique, and it is linear in $T$.

\paragraph{The action (condition K1) and the Newton constant.}
Each species is cut in its own operator:
\begin{equation}
  W(L)=-\tfrac12\sum_{\text{species}} m\,\sigma\,\operatorname{Tr} f(D/L),
\end{equation}
with multiplicity $m$ and statistics sign $\sigma$. With the Euclidean convention
$W\supset-(16\pi G)^{-1}\!\int\!\sqrt g\,R$~\cite{FrolovFursaevZelnikov1997}, the content
induces $\kappa$ when $W$ has coefficient $-\kappa\int\!\sqrt g\,R$ at order $L^s$. Then
$(16\pi G)^{-1}=\kappa L^s$, and gravity attracts, $G>0$, at every scale exactly when
$\kappa>0$.

\paragraph{The theorem.}
Let $M_f=\Gamma(s)^{-1}\int_0^\infty f(u)\,u^{s-1}du$ be the profile's moment, Eq.~\eqref{eq:moment},
normalized to the heat kernel's, so that $M_f=1$ for $f(u)=e^{-u}$. If each species' cutoff trace has
coefficient $M_f a_1$ at order $L^s$ (proved for proper-time profiles in
Appendix~\ref{app:abelian}), then
\begin{equation}\label{eq:kappa}
  \kappa=\frac{M_f\,(4\pi)^{-(s+1)}}{12}\sum_{\text{species}} m\,\sigma\,(N-6c),
\end{equation}
by the linearity of~\eqref{eq:coeff}. The sum is the net count. For $M_f>0$, $\kappa>0$
exactly when the net count is positive: the profile scales the coefficient and cannot
change its sign. $M_f>0$ holds for every proper-time profile, and for every non-increasing
profile with finite moments that vanishes at infinity and is not identically zero.

%======================================================================
\section{The expansion, proved for proper-time profiles}
\label{app:abelian}

A proper-time profile is a Laplace transform,
\begin{equation}
  f(u)=\int e^{-tu}\,d\nu(t),
\end{equation}
of a non-zero finite measure $\nu$ carried by $(0,T]$ with $\int t^{-(s+1)}d\nu<\infty$. It
is the regulated logarithm of every proper-time cutoff with an ultraviolet and an infrared
end, and it contains the heat kernel's profile, $\nu=\delta_1$.

\paragraph{Claim.} For such a profile, $\operatorname{Tr}f(D/L)$ has coefficient $M_f a_1$
at order $L^s$. Nothing is assumed beyond the heat-trace bound~\eqref{eq:heatbound}.

\paragraph{Proof.}
By Fubini, $\operatorname{Tr}f(D/L)=\int\Theta(t/L)\,d\nu(t)$: the cutoff trace is the heat
trace averaged over proper time. For $L\ge T/\tau_0$, every $t$ in the support has
$t/L\le\tau_0$, and~\eqref{eq:heatbound} at $\tau=t/L$, multiplied by $L\,t^{-(s+1)}$,
gives
\begin{equation}
  \Bigl|L^{-s}\Theta(t/L)-a_0\,L\,t^{-(s+1)}-a_1\,t^{-s}\Bigr|\le \frac{|C|\,t^{1-s}}{L}.
\end{equation}
Integrating against $\nu$, with $A=a_0\int t^{-(s+1)}d\nu$,
\begin{multline}
  \Bigl|L^{-s}\bigl(\operatorname{Tr}f(D/L)-A\,L^{s+1}\bigr)-a_1\!\int t^{-s}d\nu\Bigr|\\
  \le\frac{|C|}{L}\int t^{1-s}d\nu\;\to\;0 .
\end{multline}
The moments used are finite: on $(0,T]$, $t^{-s}\le T\,t^{-(s+1)}$ and
$t^{1-s}\le T^2\,t^{-(s+1)}$, and $\int t^{-(s+1)}d\nu$ is finite by assumption. Finally,
the Mellin transform gives $\int_0^\infty f(u)u^{s-1}du=\Gamma(s)\int t^{-s}d\nu$, so
$\int t^{-s}d\nu=M_f$. \hfill$\square$

\noindent The rest of the formal proof establishes integrability for Fubini and for each
moment, and is routine.

%======================================================================
\section{One to six, derived}
\label{app:onetosix}

A species' operator is its kinetic operator $D_0$ with a curvature term inserted. Under
the cut in the species' own operator (K1), the curvature coefficient can be read in two
ways, and the two agree for every proper-time profile.

\paragraph{The full operator.}
A constant insertion $\delta$ shifts the spectrum by $\delta$, so the heat trace of
$D_0+\delta$ is $e^{-\tau\delta}\Theta(\tau)$. Expanding the product,
\begin{equation}
  \tau^{s+1}e^{-\tau\delta}\Theta(\tau)=a_0+(a_1-\delta a_0)\,\tau+O(\tau^2),
\end{equation}
with an explicit constant in the remainder. So the shifted operator has
$a_1(D_0+\delta)=a_1(D_0)-\delta\,a_0(D_0)$: Gilkey's curvature term for a constant
insertion, derived rather than imported. By Appendix~\ref{app:abelian}, its cutoff trace
has coefficient $M_f\,(a_1-\delta a_0)$.

\paragraph{The kinetic part plus the response.}
The kinetic part is the cut in $D_0$, with coefficient $M_f a_1$. The insertion, read as the
first-order response of the cutoff trace,
\begin{equation}
  \mathcal R(L)=\frac{d}{d\delta}\operatorname{Tr}f\!\left(\frac{D_0+\delta}{L}\right)\Big|_{\delta=0}
  =-\int\frac{t}{L}\,\Theta\!\left(\frac tL\right)d\nu(t),
\end{equation}
has coefficient $-a_0\,\mathcal D/\Gamma(s+1)$ at order $L^s$, where
$\mathcal D=-\int_0^\infty f'(u)\,u^s\,du$ is the profile's derivative moment. The estimate
is the one in Appendix~\ref{app:abelian}, one power of $t$ higher. Differentiating under
the integral takes uniform bounds, which are routine.

\paragraph{The integration by parts.}
The Mellin transform gives $\mathcal D=\Gamma(s+1)\int t^{-s}d\nu$, so
$\mathcal D/\Gamma(s+1)=M_f$. For proper-time profiles, the integration by parts between the
two moments is the Gamma function's recursion, $\Gamma(s+1)=s\,\Gamma(s)$.

\paragraph{One to six.}
Write $\mathcal I=(4\pi)^{-(s+1)}\int\!\sqrt g\,R$. With the kinetic operator's
$a_1=(N/6)\,\mathcal I$ and an insertion of trace $c$, so that
$\delta a_0=c\,\mathcal I$, both readings give
\begin{equation}
  M_f\,\frac{N}{6}\,\mathcal I-\delta\,a_0\,M_f=M_f\,\frac{N-6c}{6}\,\mathcal I .
\end{equation}
The kinetic part and the insertion stand at one to six, with the same factor $M_f$, for
every proper-time profile. An insertion that varies over the background is not derived
this way; for it the coefficient is Gilkey's, as imported in Appendix~\ref{app:setup}.

\paragraph{Profiles of bounded variation.}
A profile need not be differentiable; the sharp step is the obvious case. A profile of
bounded variation can be written through its derivative measure,
$f(u)=m_+((u,\infty))-m_-((u,\infty))$, with $m_\pm$ carried by $(0,\infty)$ and of finite
$s$-th moment. For each side, the layer-cake formula gives
\begin{equation}
  \int u^s\,dm=s\int_0^\infty m((t,\infty))\,t^{s-1}dt,
\end{equation}
so $\int u^s\,d(m_+-m_-)=s\int_0^\infty f(u)\,u^{s-1}du$, with no derivative and no boundary
term. The insertion's weight relative to the kinetic part is therefore six for every such
profile with non-zero moment. What turns that weighting into the induced coefficient is the
spectral-action expansion, which the literature states for Laplace-transform and for
Schwartz profiles~\cite{vanSuijlekom2024,EcksteinIochum2018,EstradaGraciaBondiaVarilly1998}
and which is taken here as a hypothesis for them, not proved. For the sharp step it holds
beyond the leading term only as a Ces\`aro mean~\cite{EstradaGraciaBondiaVarilly1998}.

%======================================================================
\section{The counts in closed form}
\label{app:counts}

\paragraph{Conventions.}
Fermions are counted as Weyl (two-component) fermions; Dirac and Majorana counts are
derived from them, never entered separately. A species' field components $N$ are its
covariant components, one per real field, the convention its heat-kernel coefficient is
stated in. The horizon count uses a different quantity, each species' coefficient in the
horizon's entanglement entropy in units of one real scalar (Table~\ref{tab:species}),
which is not a count of physical degrees of freedom. Each gauge vector is counted with
its two Faddeev--Popov ghosts, minimally coupled scalars of opposite statistics. The
contact term is assigned to the vector, for the gauge field with its ghosts, as Kabat
states it~\cite{Kabat1995}; no other species has one. Each of the \gaugeVectors{} vectors
counts as one abelian gauge field, which is exact at one loop on a purely gravitational
background.

\paragraph{The closed form.}
With the Standard Model's gauge group and its ghosts, a content of $g$ generations, $h$
Higgs doublets and $n$ right-handed neutrinos has $\weylPerGeneration\,g+n$ Weyl fermions
of weight \weightWeyl, $\scalarsPerDoublet\,h$ real scalars of weight \weightScalar,
\gaugeVectors{} vectors of weight \weightVector{} and \smGhosts{} ghosts of weight
\weightGhost. Its net count is Eq.~\eqref{eq:content},
\begin{equation}
  \Nnet(g,h,n)=\weylPerGeneration\,g+\scalarsPerDoublet\,h+n\smGaugeSector .
\end{equation}
Every statement below holds for all non-negative integers; none rests on a search.

\paragraph{The floor.}
With one doublet, $\Nnet=\weylPerGeneration\,g\genCountZero$, which is positive exactly
when $g\ge\floorGenerations$. The values for zero to six generations are those of
Fig.~\ref{fig:generations}.

\paragraph{The least content.}
A count of exactly one needs $\weylPerGeneration\,g+\scalarsPerDoublet\,h=\countOneNeeds$.
Reducing modulo four, $-g\equiv1$, so $g\equiv3\pmod 4$. Then $g=3$ gives $h=1$, and
$g\ge7$ overshoots, so three generations and one doublet is the only solution. With a
right-handed neutrino in every generation, the condition becomes
$(\weylPerGeneration+1)\,g+\scalarsPerDoublet\,h=\countOneNeeds$, which has no solution,
since the left side is a multiple of four and the right side is not. As a cross-check the
repository also searches every content with fewer than \searchedGenerations{} generations
and fewer than \searchedDoublets{} doublets, a range that covers every possible solution;
it finds the same.

\paragraph{Right-handed neutrinos.}
Adding $n$ of them to the Standard Model gives a count of $1+n$, a kinetic total of
$\smKinetic-2n$ and an insertion total of $\smInsertion+3n$. So, with $\Nkin=\smKinetic$ and
$\Nins=\smInsertion$, the threshold is $6(|\Nkin|+2n)/(\Nins+3n)$ and the margin is
$(1+n)/(\Nins+3n)$,
which increases with $n$. The sign never changes.

\paragraph{The curvature coupling.}
Give the doublet's \smScalars{} real scalars a common coupling $\xi$, so each has curvature
trace $\xi$ and weight $1-6\xi$. The Standard Model's count becomes
$\smRestWeight+\smScalars\,(1-6\xi)=1-\smCouplingSlope\,\xi$, positive exactly when
$\xi<\smCouplingBound$. In general, a content whose other species contribute $r$ and whose
$k$ scalars share the coupling is positive exactly when $\xi<(r+k)/(6k)$, which gives
\nuTwoCouplingBound{} and \nuThreeCouplingBound{} with two and three right-handed
neutrinos. A single scalar coupled on its own faces a different bound; it is the common
coupling that gives \smCouplingBound. At the conformal coupling $\xi=1/6$ the count is
\smConformalCount, \nuTwoConformalCount{} and \nuThreeConformalCount{} with none, two
and three right-handed neutrinos.

\paragraph{Gauge bosons.}
A content of $k_s$ minimally coupled real scalars, $k_W$ Weyl fermions and $v$ gauge
vectors, each with its two ghosts, has $\Nnet=k_s+k_W-\matterPerBoson\,v$. With the
Standard Model's \smScalars, \smWeyl{} and \gaugeVectors{} this is \smNet, so $b$ further
vectors and $k$ further scalars or Weyl fermions give $1+k-\matterPerBoson\,b$, which is
positive exactly when $k\ge\matterPerBoson\,b$, since both sides are integers.

\paragraph{The horizon identity.}
Over any content built from the four species, the horizon count, entanglement
coefficients plus contact terms, differs from the induced count by
\begin{equation}
  (\text{entanglement}+\text{contact})-\Nnet=n_{\text{ghost}}-2\,n_{\text{vector}} .
\end{equation}
So the two agree exactly when every vector carries its two ghosts, which is condition K2.
For the Standard Model they read \smEntanglement{} and \smContact, a net of \smNet. That the
horizon count \emph{should} equal the induced one is the proportionality between horizon
entropy and the induced coupling~\cite{LarsenWilczek1996,FrolovFursaevZelnikov1997}, which
is cited here, not proved.

%======================================================================
\section{The magnitude}
\label{app:magnitude}

% The Lean table is set here, a page early, so that it floats ahead of the references.
\begin{table*}
\caption{The results that cite Lean: each result's key in the repository's
\texttt{CLAIMS.md}, the section in which it appears, and the Lean statements it cites.}
\label{tab:lean}
\centering
\small
\input{../../build/latex/tables/lean.tex}
\end{table*}

In four dimensions, Eq.~\eqref{eq:kappa} gives $(16\pi G)^{-1}=M_f\,\Nnet\,L/(192\pi^2)$, so
a content of net count $\Nnet$ induces $G=M_P^{-2}$ exactly when
\begin{equation}
  \Lambda=M_P\sqrt{\frac{12\pi}{M_f\,\Nnet}} .
\end{equation}
The normalized moment is $1/\Gamma(s+1)$ for the sharp step $\theta(1-u)$, $1$ for the
heat kernel's profile, and $\Gamma(s/2)/(2\Gamma(s))$ for the Gaussian $e^{-u^2}$. In four
dimensions these are $1$, $1$ and $\sqrt\pi/2\approx\momentGaussian$. For the Standard
Model this gives \smCutoffLow{} to \smCutoffHigh\,$M_P$ (Table~\ref{tab:magnitude}). The
heat kernel's profile is a proper-time profile, so its row rests on
Appendix~\ref{app:abelian}; the Gaussian's rests on the imported expansion for Schwartz
profiles; and the sharp step's rests on an expansion that holds beyond the leading term
only as a Ces\`aro mean (Appendix~\ref{app:onetosix}). The integration by parts, and these
moments, are also checked numerically in the repository against the closed forms, at
$s=1$ and $s=2$, with the quadrature's tolerance checked again in a second software environment with
independently installed library versions. The magnitude moves with the profile; the sign does not.

%======================================================================
\section{The formal statements}
\label{app:formal}

Every result in Appendices~\ref{app:setup}--\ref{app:magnitude} is proved in Lean~4
against Mathlib, in the repository's Lean library, under the checks described in
Sec.~\ref{sec:lean}; imported results enter only as hypotheses the theorems take.
Table~\ref{tab:lean} lists every result of the main text that cites Lean: its key in
the repository's \texttt{CLAIMS.md}, where its wording and evidence are given; the
section of this paper in which it appears; and the Lean statements it cites.

\bibliography{../../build/latex/refs}

\end{document}
